\chapter{Generalisasi, Regularisasi, dan Pelatihan}

Sebuah tim keamanan membangun pendeteksi intrusi dari dua besaran mentah: banyak
permintaan per detik dan ukuran muatan rata-rata. Dari dua besaran itu mereka
menurunkan dua puluh satu ciri dengan mengalikan dan memangkatkan keduanya sampai
derajat lima, lalu melatih pengklasifikasi logistik pada enam puluh koneksi berlabel.
Pada data latihan, pendeteksi itu salah pada nol dari enam puluh koneksi. Pada dua
puluh ribu koneksi baru yang tidak pernah dilihatnya, ia salah pada tiga puluh dua
persen. Model yang sama, bobot yang sama, dua angka yang berbeda jauh.

Angka nol itu bukan kebetulan dan bukan prestasi. Dengan dua puluh satu ciri dan enam
puluh contoh, sebuah bidang pemisah selalu dapat ditemukan yang memotong setiap contoh
latih dengan benar, termasuk contoh yang labelnya keliru. Yang dipelajari model itu
bukan bentuk pemisah yang benar, melainkan posisi enam puluh titik itu. Begitu titiknya
berganti, bidangnya salah tempat.

Bab 25 mengajari kita menghitung gradien dan menurunkannya sampai galat latih kecil;
Bab 20 mengajari kita membaca satu angka beserta ketidakpastiannya. Tidak satu pun
menjawab pertanyaan yang menentukan di sini: besaran apa yang harus dikecilkan supaya
angka di lapangan ikut mengecil, dan bagaimana memilih setelan latihan tanpa menipu diri
sendiri.

\begin{untukanak}
Seorang murid diberi enam puluh soal latihan menjelang ujian. Ia belajar dengan cara
menghafal: setiap soal beserta jawabannya, sampai huruf per huruf. Pada soal latihan ia
selalu benar, seratus dari seratus, dan ia pulang dengan bangga. Pada ujian, soalnya
berbeda: angkanya lain, kalimatnya lain, urutannya lain. Nilainya jatuh, dan gurunya
bingung karena murid ini tampak paling rajin.

Murid kedua belajar lebih lambat. Ia tidak menghafal jawaban, ia mencari polanya: kapan
harus dibagi, kapan harus dikali, kapan jawabannya masuk akal. Pada soal latihan ia
salah pada beberapa nomor, dan ia tidak bangga. Pada ujian ia justru lebih baik.

Sekarang bagian yang membuat penasaran. Andaikan murid pertama diberi dua ribu soal
latihan, bukan enam puluh. Apakah menghafal masih buruk? Ternyata tidak selalu: kalau
soalnya cukup banyak dan bervariasi, menghafal menjadi cara belajar yang mahal tetapi
bekerja. Jadi masalahnya bukan menghafal, melainkan menghafal dari contoh yang sedikit.
Di bab ini kita akan tahu mengapa jumlah contoh dan kerumitan cara belajar harus
seimbang, mengapa menambah kerumitan justru memperburuk nilai ujian, dan bagaimana
mengukur keseimbangan itu tanpa harus menunggu ujian.
\end{untukanak}

\section{Dua galat dan satu jurang}
\tingkat{dasar}

Setiap model yang dilatih punya dua angka, dan keduanya dihitung dengan cara sama
persis: proporsi kesalahan pada sekumpulan contoh. Bedanya hanya himpunan mana yang
dipakai. Kalimat sehari-harinya: nilai ujian di rumah berbeda dari nilai ujian di
sekolah, walaupun soalnya sama jenisnya.

\kamus{generalisasi}{generalization}{kemampuan mengerjakan soal baru, bukan soal yang
sudah dihafal}

\begin{definisi}[Galat latih, galat populasi, dan jurang]
\index{generalisasi}
Misalkan contoh $(x_i,y_i)$ diambil dari distribusi $\mathcal{D}$ yang tidak diketahui.
Untuk hipotesis $h$ dan galat 0-1, \emph{galat populasi} adalah
$R(h)=\mathbb{P}_{(x,y)\sim\mathcal{D}}\bigl(h(x)\neq y\bigr)$, sedangkan \emph{galat
latih} pada $n$ contoh adalah
$\hat R_n(h)=\frac1n\sum_{i=1}^n \mathbf{1}\{h(x_i)\neq y_i\}$.
Selisih $R(h)-\hat R_n(h)$ disebut \emph{jurang generalisasi} untuk $h$ itu.
\end{definisi}

Galat latih dapat dihitung; galat populasi tidak, sebab distribusinya tidak diketahui.
Yang dapat kita lakukan hanyalah menaksirnya pada himpunan lain yang tidak dipakai
melatih, dan penaksir itu sendiri bergoyang seperti setiap proporsi pada Bab 20. Jadi
setiap angka "galat validasi" dalam bab ini adalah taksiran, bukan kebenaran.

Sekarang dua penyakit yang berlawanan. Kalimat sehari-harinya: model yang terlalu
sederhana tidak mampu menangkap polanya, sedangkan model yang terlalu rumit menghafal
kebetulan.

\begin{definisi}[Underfitting dan overfitting]
\index{underfitting}
\index{overfitting}
Sebuah model \emph{underfitting} bila galat latih dan galat populasi keduanya tinggi:
modelnya terlalu kaku untuk bentuk yang sebenarnya. Sebuah model \emph{overfitting}
bila galat latih rendah tetapi jurang generalisasinya besar: modelnya menyesuaikan diri
pada kebetulan data latih, bukan pada bentuk yang sebenarnya.
\end{definisi}

Kedua penyakit itu tidak simetris dalam praktik. Underfitting terlihat jelas: galat latih
tinggi, dan menambah ciri menolong. Overfitting menipu, sebab satu-satunya angka yang
tersedia dengan mudah, yaitu galat latih, justru tampak sempurna. Model pada pembuka bab
adalah contoh paling murni: bobot terbesarnya $406{,}204$, tanda bahwa skornya bergantung
pada perbedaan kecil di antara titik, bukan pada skala yang masuk akal.

Yang mengatur keduanya adalah \emph{kapasitas}: seberapa kaya himpunan fungsi yang boleh
dipilih model. Untuk pengklasifikasi linear atas $d$ ciri, kapasitas naik dengan $d$;
untuk jaringan saraf, dengan jumlah unit, kedalaman, dan jumlah langkah latihan.
Kapasitas bukan sekadar cacah parameter: model besar yang diregularisasi kuat dapat
berperilaku seperti model kecil.

\begin{bacakeras}
Dua suara, bergantian.

``Seratus dari seratus. Berarti ia paham, bukan?''

``Belum tentu. Seratus dari seratus hanya berarti tidak ada soal yang baru.''

``Jadi menghafal itu salah?''

``Menghafal bukan salah. Menghafal dari sedikit contoh yang salah.''

``Bedanya di mana?''

``Di soal yang tidak ada di catatan. Hafalan menjawab yang sudah pernah dilihat.
Pemahaman menjawab yang belum pernah dilihat.''

``Dan kalau soalnya dua ribu, bukan enam puluh?''

``Maka hafalan mulai menolong. Bukan karena hafalan berubah menjadi pemahaman, tetapi
karena kebetulan yang tersimpan semakin banyak yang terpakai lagi.''

Ucapkan sekali lagi, pelan: hafalan menjawab yang sudah ada, pemahaman menjawab yang
belum ada. Ulangi baris itu sampai terasa wajar di mulut: sudah ada, belum ada.
\end{bacakeras}

Jurang generalisasi menyusut ketika contoh bertambah. Pada eksperimen kapasitas dengan
anggaran langkah tetap, jurangnya $0{,}1338$ pada $n=60$, $0{,}0274$ pada $n=150$,
$0{,}0193$ pada $n=600$, dan $0{,}0097$ pada $n=2400$, sedangkan galat latihnya hampir
tidak bergerak di kisaran $0{,}17$ sampai $0{,}21$.

\section{Dekomposisi bias-variansi}
\tingkat{menengah}

Pertanyaan yang lebih tajam: dari mana jurang itu berasal. Kalimat sehari-harinya:
seorang penembak dapat meleset karena bidikannya bergeser ke satu arah, atau karena
tangannya gemetar, dan dua sebab itu menuntut obat yang berbeda.

\kamus{bias-variansi}{bias-variance}{meleset sistematis lawan gemetar; keduanya
menyumbang galat, dan menekan yang satu kadang menaikkan yang lain}

\begin{teorema}[Dekomposisi bias-variansi]
\index{bias-variansi}
Misalkan $y=f(x)+\varepsilon$ dengan $\mathbb{E}[\varepsilon]=0$ dan
$\operatorname{Var}(\varepsilon)=\sigma^2$, dan misalkan $\hat f$ dilatih pada
himpunan latih acak $D$ yang bebas dari $\varepsilon$. Untuk setiap $x$ tetap,
\[
\mathbb{E}_{D,\varepsilon}\bigl[(y-\hat f(x))^2\bigr]
=\sigma^2
+\bigl(\mathbb{E}_D\hat f(x)-f(x)\bigr)^2
+\operatorname{Var}_D\bigl(\hat f(x)\bigr).
\]
\end{teorema}

\begin{proof}[Bukti]
Tulis $y-\hat f(x)=\varepsilon+\bigl(f(x)-\mathbb{E}_D\hat f(x)\bigr)
+\bigl(\mathbb{E}_D\hat f(x)-\hat f(x)\bigr)$, yaitu tiga suku: derau yang tidak dapat
dikurangi, bias, dan simpangan dari rata-rata. Kuadratkan dan ambil harapannya. Suku
silang yang memuat $\varepsilon$ lenyap karena $\mathbb{E}[\varepsilon]=0$ dan
$\varepsilon$ bebas dari $D$; suku silang antara bias dan simpangan lenyap karena
$\mathbb{E}_D[\hat f(x)-\mathbb{E}_D\hat f(x)]=0$. Yang tersisa tepat tiga suku pada
pernyataan teorema.
\end{proof}

Suku pertama tidak dapat dikurangi model apa pun; ia milik datanya, disebut galat tak
tereduksikan, dan pada data pembuka bab nilainya sekitar $0{,}21$ karena labelnya acak.
Suku kedua dan ketiga dapat diperdagangkan: menambah kapasitas menurunkan bias dan
menaikkan variansi, mengurangi kapasitas melakukan sebaliknya. Regularisasi adalah nama
untuk perdagangan itu.

Contoh paling murah yang memuat kedua suku sekaligus adalah penaksir menyusut. Untuk
mengukur $\mu$ dari $n$ pengamatan berderau $\sigma^2$, orang biasa memakai
$\bar X_n$, yang tak bias. Kita boleh memakai $c\bar X_n$ untuk $c\in(0,1)$, dan
penaksir itu berbias dengan sengaja. Galat kuadrat rata-ratanya
\[
\operatorname{MSE}(c)=(c-1)^2\mu^2+c^2\frac{\sigma^2}{n},
\]
diminimalkan di $c^\ast=\mu^2/(\mu^2+\sigma^2/n)$. Untuk $\mu=2$, $\sigma=1$, dan
$n=25$ diperoleh $c^\ast=0{,}990099$: penaksir terbaik justru sedikit menyusut.

\begin{center}
\small
\begin{tabular}{rrrr}
\hline
$c$ & bias$^2$ & variansi & MSE \\
\hline
0,70 & 0,360000 & 0,019600 & 0,379600 \\
0,90 & 0,040000 & 0,032400 & 0,072400 \\
0,99 & 0,000400 & 0,039204 & 0,039604 \\
1,00 & 0,000000 & 0,040000 & 0,040000 \\
1,10 & 0,040000 & 0,048400 & 0,088400 \\
1,30 & 0,360000 & 0,067600 & 0,427600 \\
\hline
\end{tabular}
\end{center}

Tiga hal perlu dibaca dari tabel itu. Keuntungan menyusut kecil: MSE turun dari
$0{,}040000$ ke $0{,}039604$, sekitar satu persen. Kerugian menyimpang ke arah yang
salah jauh lebih besar: $c=1{,}3$ memberi $0{,}4276$, sedangkan $c=0{,}7$ hanya
$0{,}3796$. Dan minimumnya landai, sehingga memilih $c$ yang sedikit keliru bukan
bencana.

Teorema itu tidak memberi rumus untuk $\hat f$ yang rumit seperti jaringan saraf.
Dekomposisinya berlaku formal, tetapi kedua sukunya tidak dapat dihitung dari data saja;
yang teramati hanyalah jumlahnya, yaitu galat validasi. Karena itu sisa bab ini bekerja
dengan angka teramati, bukan dengan bias dan variansi yang sebenarnya.

\section{Regularisasi L2 dan L1}
\tingkat{menengah}

Cara paling langsung menekan variansi adalah melarang bobot tumbuh besar. Kalimat
sehari-harinya: kita menambahkan denda ke dalam sasaran, sehingga model harus membayar
untuk setiap bobot besar, dan hanya mau membayar bila keuntungannya di data latih sepadan.

\kamus{regularisasi}{regularization}{denda yang dibayar model untuk bobot besar, supaya
ia tidak menghafal data latih}

\begin{definisi}[Sasaran ternormalisasi]
\index{regularisasi}
Untuk sasaran asal $J(\mathbf{w})$ dan penalti $\Omega$, sasaran yang diminimalkan
adalah $\tilde J(\mathbf{w})=J(\mathbf{w})+\lambda\Omega(\mathbf{w})$ dengan
$\lambda\ge0$ sebagai hiperparameter. Untuk penalti L2, $\Omega=\frac12\|\mathbf{w}\|_2^2$;
untuk penalti L1, $\Omega=\|\mathbf{w}\|_1=\sum_j|w_j|$.
\end{definisi}

Untuk L2, gradien sasarannya $\nabla J+\lambda\mathbf{w}$, jadi satu langkah penurunan
gradien dengan laju $\eta$ berbentuk
\[
\mathbf{w}\leftarrow\mathbf{w}-\eta\bigl(\nabla J+\lambda\mathbf{w}\bigr)
=(1-\eta\lambda)\mathbf{w}-\eta\nabla J .
\]
Faktor $(1-\eta\lambda)$ itulah asal nama \emph{weight decay}: pada setiap langkah,
sebelum gradien dipakai, bobot digandakan dengan bilangan yang sedikit kurang dari satu.
Kalau gradiennya nol, bobotnya meluruh menuju nol secara geometris dengan waktu paruh
$\ln 2/(\eta\lambda)$ langkah.

\begin{akibat}[Peluruhan bobot sama dengan penalti L2]
\index{weight decay}
Untuk laju $\eta$ dan penalti L2 dengan koefisien $\lambda$, penurunan gradien atas
sasaran ternormalisasi identik dengan penurunan gradien atas sasaran asal yang diikuti
penggandaan bobot dengan $1-\eta\lambda$. Dengan demikian $\lambda$ dan $\eta$ hanya
muncul sebagai hasil kali $\eta\lambda$ pada langkah peluruhan, sehingga menaikkan laju
dua kali dan menurunkan $\lambda$ dua kali memberi lintasan peluruhan yang sama.
\end{akibat}

Penalti L2 juga punya tafsiran Bayesian: ia setara prior Gaussian pada bobot, sehingga
penyelesaiannya taksiran MAP, bukan kemungkinan maksimum. Itu menjelaskan mengapa L2
percaya bobot kecil lebih mungkin, dan mengapa ia tidak pernah memaksa satu bobot menjadi
nol.

L1 berperilaku lain, dan seluruh bedanya terletak pada bentuk penaltinya di dekat nol.
Kalimat sehari-harinya: penalti kuadratik menarik semakin lemah ketika bobotnya sudah
kecil, sedangkan penalti nilai mutlak menarik dengan kekuatan tetap sampai bobotnya
benar-benar nol.

\begin{lemma}[Ambang lunak]
\index{regularisasi}
Untuk $a\in\mathbb{R}$ dan $t>0$, peminimum $u\mapsto\frac12(u-a)^2+t|u|$ adalah
\[
u^\ast=\operatorname{sign}(a)\max\{|a|-t,0\},
\]
sehingga $u^\ast=0$ tepat ketika $|a|\le t$.
\end{lemma}

Untuk $u>0$ fungsinya $\frac12(u-a)^2+tu$ dengan turunan $u-a+t$, jadi titik
kritisnya $u=a-t$, sah hanya bila $a>t$. Untuk $u<0$ turunannya $u-a-t$ dengan titik
kritis $u=a+t$, sah hanya bila $a<-t$. Pada $u=0$ subgradiennya selang $[-a-t,-a+t]$,
yang memuat nol tepat ketika $|a|\le t$; dalam hal itu nol peminimum karena fungsinya
konveks. Ketiga kasus bergabung menjadi rumus di atas.

Lemma itu menjelaskan seluruh perilaku L1. Dengan pendekatan kuadratik di sekitar
penyelesaian tanpa penalti $w^\ast$ dan kelengkungan $H$, sasarannya kira-kira
$\frac H2(w-w^\ast)^2+\lambda|w|$, sehingga
$w=\operatorname{sign}(w^\ast)\max\{|w^\ast|-\lambda/H,0\}$: setiap bobot di bawah
$\lambda/H$ menjadi nol tepat, bukan sekadar kecil. Untuk L2 penyelesaiannya
$w=w^\ast/(1+\lambda/H)$, yang tidak pernah nol kecuali $w^\ast$ sendiri nol.

\begin{proposisi}[Mengapa L1 menghasilkan nol sejati]
\index{regularisasi}
Misalkan $J$ dapat diturunkan dan konveks di sekitar nol. Titik $\mathbf{w}=\mathbf{0}$
adalah peminimum $J(\mathbf{w})+\lambda\|\mathbf{w}\|_1$ jika dan hanya jika
$|\partial J/\partial w_j|\le\lambda$ untuk setiap $j$. Untuk penalti L2, syarat yang
sama adalah $\partial J/\partial w_j=0$ untuk setiap $j$.
\end{proposisi}

\begin{proof}[Bukti]
Subgradien $\|\cdot\|_1$ di nol adalah kotak $[-1,1]^d$, jadi kondisi
$0\in\nabla J(\mathbf{0})+\lambda[-1,1]^d$ setara dengan $|(\nabla J)_j|\le\lambda$
untuk semua $j$; untuk fungsi konveks, kondisi itu perlu dan cukup. Untuk penalti L2,
gradiennya $\nabla J+\lambda\mathbf{w}$ yang di $\mathbf{0}$ sama dengan $\nabla J$, jadi
syaratnya $\nabla J(\mathbf{0})=\mathbf{0}$.
\end{proof}

Akibat praktisnya penting untuk pembacaan hasil. L1 memberi daerah datar selebar
$2\lambda$ di sekitar nol, sehingga bobot yang "hampir tidak berguna" dimatikan
sepenuhnya; L2 hanya mengecilkannya. Karena itu L1 dipakai sebagai pemilih ciri, dan
karena itu pula jumlah bobot nol pada L1 bukan ukuran kepentingan yang halus: ia
melompat, tidak bergeser perlahan.

Sekarang angkanya. Pengklasifikasi logistik yang sama, data yang sama, dua puluh satu
ciri pada enam puluh contoh latih, semuanya dihitung dengan skrip pada bagian lab.

\begin{center}
\small
\begin{tabular}{rrrr}
\hline
$\lambda$ & $\max_j|w_j|$ & galat latih & galat validasi \\
\hline
0       & 406,204 & 0,0000 & 0,3262 \\
0,001   &   3,056 & 0,1000 & 0,2927 \\
0,003   &   1,814 & 0,1000 & 0,2688 \\
0,010   &   1,018 & 0,1167 & 0,2475 \\
0,030   &   0,631 & 0,1167 & 0,2433 \\
0,100   &   0,399 & 0,1500 & 0,2406 \\
0,300   &   0,268 & 0,2333 & 0,2452 \\
1,000   &   0,142 & 0,2500 & 0,2695 \\
3,000   &   0,062 & 0,3333 & 0,3205 \\
10,000  &   0,022 & 0,3333 & 0,3395 \\
100,000 &   0,003 & 0,3333 & 0,3395 \\
\hline
\end{tabular}
\end{center}

\begin{figure}[htbp]
\centering
\begin{tikzpicture}
\begin{axis}[
  width=12.2cm, height=5.6cm,
  xmin=-0.4, xmax=10.4, ymin=0, ymax=0.40,
  axis lines=left,
  xtick={0,3,5,7,10},
  xticklabels={$0$,$10^{-2}$,$0{,}1$,$1$,$100$},
  ytick={0,0.1,0.2,0.3,0.4},
  yticklabels={$0$,$0{,}10$,$0{,}20$,$0{,}30$,$0{,}40$},
  tick label style={font=\scriptsize},
  xlabel={kekuatan regularisasi $\lambda$ (kapasitas menurun ke kanan)},
  ylabel={galat},
  xlabel style={font=\small},
  ylabel style={font=\small},
  legend style={font=\scriptsize,draw=none,at={(0.03,0.97)},anchor=north west},
]
\addplot[very thick,layerindigo,mark=*,mark size=1.5pt] coordinates {
  (0,0.0000) (1,0.1000) (2,0.1000) (3,0.1167) (4,0.1167) (5,0.1500)
  (6,0.2333) (7,0.2500) (8,0.3333) (9,0.3333) (10,0.3333)};
\addlegendentry{galat latih}
\addplot[very thick,accent,mark=square*,mark size=1.5pt] coordinates {
  (0,0.3262) (1,0.2927) (2,0.2688) (3,0.2475) (4,0.2433) (5,0.2406)
  (6,0.2452) (7,0.2695) (8,0.3205) (9,0.3395) (10,0.3395)};
\addlegendentry{galat validasi}
\addplot[only marks,mark=o,mark size=3.4pt,accent] coordinates {(5,0.2406)};
\draw[densely dashed,layergrey] (axis cs:5,0)--(axis cs:5,0.2406);
\node[accent,font=\scriptsize,anchor=north] at (axis cs:4.55,0.238)
  {minimum $0{,}2406$};
\end{axis}
\end{tikzpicture}
\caption{Galat latih dan galat validasi terhadap kekuatan regularisasi $\lambda$; arah ke
kanan berarti kapasitas menurun, sebab bobot ditahan makin ketat. Kesebelas titik adalah
baris pada tabel sebelumnya. Galat latih naik monoton dari $0{,}0000$ ke $0{,}3333$,
sedangkan galat validasi membentuk lembah: $0{,}3262$ pada $\lambda=0$, minimum
$0{,}2406$ pada $\lambda=0{,}1$, lalu $0{,}3395$ pada $\lambda=10$ ke atas. Dasar
lembahnya lebar, dari $0{,}2433$ pada $\lambda=0{,}03$ ke $0{,}2452$ pada $\lambda=0{,}3$.}
\label{fig:lembah-kapasitas}
\end{figure}

Gambar~\ref{fig:lembah-kapasitas} memplot kolom ketiga dan keempat tabel itu sebagai
satu lembah. Bacalah tabel itu sebagai tiga daerah, bukan sebelas baris. Tanpa
regularisasi galat latihnya nol dan galat validasinya $0{,}3262$, dengan bobot sampai
ratusan karena tidak ada yang menahannya. Di daerah $\lambda$ antara $0{,}001$ dan
$0{,}1$, galat latih naik ke $0{,}1500$ sementara galat validasi turun ke $0{,}2406$.
Di daerah $\lambda\ge3$ keduanya buruk: $0{,}3333$ dan $0{,}3395$.

Baris terakhir pantas ditekankan. Dengan $\lambda=10$ galat validasinya $0{,}3395$,
\emph{lebih buruk} daripada $0{,}3262$ milik model tanpa regularisasi. Jadi
regularisasi bukan obat yang selalu menolong; ia menolong pada dosis tertentu dan merusak
di luar itu. Pada $\lambda=10$ model itu praktis menebak selalu kelas mayoritas, yang
galatnya $0{,}3419$: bobotnya begitu kecil sehingga skornya hampir seluruhnya ditentukan
bias.

Penalti L1 memberi gambaran serupa dengan bentuk berbeda. Pada $\lambda=0{,}05$ galat
validasinya $0{,}2334$ dengan sebelas dari dua puluh satu bobot hidup; pada
$\lambda=0{,}2$ tinggal satu bobot dengan galat $0{,}3395$. Cacah bobot hidup tidak turun
monoton: $\lambda=0{,}005$ menyisakan sebelas, $\lambda=0{,}01$ menyisakan empat belas.
Angka itu stabil terhadap empat kali lipat iterasi, jadi ia bukan kegagalan pengoptimum;
lintasan L1 untuk loss logistik memang tidak menjamin cacah nol yang monoton, dan itu
alasan tambahan untuk tidak membacanya sebagai ukuran kepentingan.

Perbandingan terakhir: L1 mencapai $0{,}2334$ dan L2 $0{,}2406$. Selisih $0{,}0072$
lebih kecil daripada goyangan antarpartisi, sekitar $0{,}02$, sehingga data ini tidak
cukup untuk menyimpulkan L1 lebih baik; yang dapat disimpulkan hanyalah bahwa keduanya
memperbaiki model tanpa regularisasi.

\begin{untukanak}
Bayangkan sebuah kelas yang harus menebak berat badan temannya. Gurunya memberi dua
cara. Cara pertama: setiap tebakan yang terlalu jauh dari jawaban dikenai denda yang
makin besar makin jauh. Cara kedua: setiap tebakan yang tidak persis sama dengan
jawaban dikenai denda tetap, tidak peduli bedanya sedikit atau banyak.

Cara pertama membuat semua tebakan bergeser sedikit ke arah tengah, tetapi tidak ada
yang menjadi nol. Cara kedua membuat tebakan yang sudah cukup dekat langsung dibulatkan
menjadi nol, karena dendanya tetap dan tidak sepadan dengan perbaikan yang mungkin
didapat. Itulah bedanya dua denda yang tampak mirip: yang satu mengecilkan, yang lain
mematikan.

Pertanyaannya: kalau kita punya dua puluh satu ciri dan hanya enam puluh contoh, denda
seperti apa yang membuat model berhenti menghafal? Jawabannya bukan "denda besar",
sebab denda yang terlalu besar mematikan semuanya, termasuk yang benar. Di bab ini kita
akan tahu mengapa ada dosis, mengapa dosis itu harus dicari, dan mengapa dosis yang
dicari pada data yang salah akan menipu kita dengan rapi.
\end{untukanak}

\section{Berhenti lebih awal, dropout, dan augmentasi}
\tingkat{menengah}

Tiga cara lain menekan variansi tanpa menyentuh rumus sasaran. Kalimat sehari-harinya:
kita dapat berhenti belajar sebelum kebetulan mulai dihafal, kita dapat mematikan
sebagian ingatan secara acak, dan kita dapat memperbanyak contoh dengan mengubah contoh
yang ada.

Berhenti lebih awal adalah bentuk regularisasi yang paling murah, sebab ia tidak
menambah satu baris pun pada sasaran: ia hanya menghentikan lintasan. Selama latihan,
galat latih hampir selalu turun, tetapi galat validasi turun lalu naik. Titik terendah
galat validasi itulah yang diambil, dan jumlah langkah menjadi hiperparameter yang
dicari.

\begin{center}
\small
\begin{tabular}{rrrr}
\hline
epoch & galat latih & galat validasi & $\max_j|w_j|$ \\
\hline
1    & 0,1333 & 0,2481 & 0,15 \\
5    & 0,1333 & 0,2429 & 0,38 \\
10   & 0,1333 & 0,2432 & 0,48 \\
50   & 0,1167 & 0,2454 & 0,77 \\
200  & 0,1167 & 0,2541 & 1,38 \\
800  & 0,1000 & 0,2879 & 2,85 \\
2400 & 0,0833 & 0,3109 & 4,81 \\
\hline
\end{tabular}
\end{center}

\begin{figure}[htbp]
\centering
\begin{tikzpicture}
\begin{axis}[
  width=11.8cm, height=5.6cm,
  xmode=log,
  xmin=1, xmax=3000, ymin=0.05, ymax=0.35,
  axis lines=left,
  xtick={1,10,100,1000},
  xticklabels={$1$,$10$,$100$,$1000$},
  ytick={0.10,0.15,0.20,0.25,0.30},
  yticklabels={$0{,}10$,$0{,}15$,$0{,}20$,$0{,}25$,$0{,}30$},
  tick label style={font=\scriptsize},
  xlabel={epoch pelatihan (skala log)},
  ylabel={galat},
  xlabel style={font=\small},
  ylabel style={font=\small},
  legend style={font=\scriptsize,draw=none,at={(0.03,0.97)},anchor=north west},
]
\addplot[very thick,layerteal,mark=*,mark size=1.5pt] coordinates {
  (1,0.1333) (5,0.1333) (10,0.1333) (50,0.1167) (200,0.1167) (800,0.1000)
  (2400,0.0833)};
\addlegendentry{galat latih}
\addplot[very thick,accent,mark=square*,mark size=1.5pt] coordinates {
  (1,0.2481) (5,0.2429) (10,0.2432) (50,0.2454) (200,0.2541) (800,0.2879)
  (2400,0.3109)};
\addlegendentry{galat validasi}
\draw[densely dashed,layergrey] (axis cs:8,0.05)--(axis cs:8,0.2427);
\addplot[only marks,mark=o,mark size=3.4pt,accent] coordinates {(8,0.2427)};
\node[accent,font=\scriptsize,anchor=north west] at (axis cs:8.6,0.2325)
  {epoch 8: $0{,}2427$};
\end{axis}
\end{tikzpicture}
\caption{Kedua galat sepanjang satu lintasan pelatihan, sumbu mendatar skala log. Galat
latih turun tanpa pernah naik, dari $0{,}1333$ pada epoch 1 ke $0{,}0833$ pada epoch
2400. Galat validasi turun sampai $0{,}2427$ pada epoch 8 lalu naik ke $0{,}3109$. Garis
putus-putus menandai titik henti terbaik: berhenti di situ memberi galat $0{,}0682$ lebih
rendah daripada melatih sampai epoch 2400, dan kurvanya mulai berpisah tepat sesudah
titik itu.}
\label{fig:dua-kurva}
\end{figure}

Gambar~\ref{fig:dua-kurva} menggambar kedua kolom itu. Galat latih turun sepanjang
lintasan, dari $0{,}1333$ ke $0{,}0833$, tanpa pernah naik, sedangkan galat validasi
turun sampai epoch 8 dengan nilai terbaik $0{,}2427$ lalu naik sampai $0{,}3109$ pada
epoch 2400. Kolom terakhir menjelaskan mengapa: yang tumbuh sepanjang latihan adalah
besarnya bobot, dan itulah variansi yang sedang kita bayar.

Berhenti lebih awal bukan sekadar taktik. Untuk sasaran kuadratik dan laju kecil, berhenti
pada langkah $\tau$ hampir setara dengan penalti L2 berkoefisien sekitar
$1/(\tau\eta)$; itu perkiraan, bukan identitas tepat. Karena itu jumlah langkah adalah
hiperparameter kapasitas, bukan urusan "melatih sampai selesai".

\begin{resep}
\begin{enumerate}
  \item Pisahkan tiga himpunan sebelum satu langkah latihan pun dijalankan: latih,
        validasi, uji. Himpunan uji disegel.
  \item Latih pada himpunan latih saja; setiap beberapa epoch hitung galat pada
        himpunan validasi, dan catat pasangan (epoch, galat validasi, galat latih).
  \item Ambil epoch dengan galat validasi terkecil, lalu kembalikan bobot ke titik itu;
        jangan memakai bobot epoch terakhir.
  \item Baru sekarang buka himpunan uji, hitung galatnya satu kali, dan laporkan
        bersama epoch pilihan serta lebar kurva validasinya.
\end{enumerate}
Peringatan urutan: langkah 1 harus selesai sebelum langkah 2, dan langkah 4 tidak boleh
dipakai menggeser pilihan pada langkah 3. Bila himpunan uji dilihat sewaktu memilih
epoch, ia bukan lagi uji melainkan validasi kedua, dan angkanya berhenti menjadi
taksiran yang jujur.
\end{resep}

Dropout mengambil jalan lain: selama latihan setiap unit dimatikan secara acak dengan
peluang $1-p$, dan yang hidup dibagi $p$ supaya harapannya tidak bergeser. Kalimat
sehari-harinya: setiap langkah sebagian anggota tim tidak masuk, sehingga tidak ada
anggota yang boleh diandalkan sendirian.

\begin{lemma}[Harapan dan variansi topeng dropout]
\index{dropout}
Misalkan $m=B/p$ dengan $B\sim\mathrm{Bernoulli}(p)$ dan $0<p\le1$. Maka
$\mathbb{E}[m]=1$ dan $\operatorname{Var}(m)=(1-p)/p$. Untuk unit dengan masukan
$z_1,\dots,z_k$ dan topeng bebas $m_1,\dots,m_k$, keluarannya
$s=\sum_j m_jw_jz_j$ memenuhi $\mathbb{E}[s]=\sum_jw_jz_j$ dan
$\operatorname{Var}(s)=\frac{1-p}{p}\sum_j(w_jz_j)^2$.
\end{lemma}

\begin{proof}[Bukti]
$\mathbb{E}[m]=p/p=1$ dan $\operatorname{Var}(m)=\mathbb{E}[m^2]-1=p/p^2-1=(1-p)/p$.
Untuk $s$, kelinearan harapan memberi suku pertama, dan kebebasan topeng memberi
variansi jumlah, yaitu $\sum_j(w_jz_j)^2\operatorname{Var}(m_j)$.
\end{proof}

Rumus variansi itu menjelaskan dua hal yang sering dilupakan. Derau dropout sebanding
dengan jumlah kuadrat kontribusi, jadi unit dengan bobot besar menerima derau relatif
lebih besar, dan dropout searah dengan penalti L2. Dan derau itu hilang sama sekali saat
pengujian, sehingga model yang dipakai berbeda dari model yang dilatih; pada jaringan
sangat kecil harga itu bisa melebihi manfaatnya.

Augmentasi data memperbanyak contoh dengan mengubah contoh yang ada, misalnya menggeser
atau menambahkan derau kecil, asalkan labelnya tidak berubah. Untuk model linear,
hubungannya dengan L2 dapat dihitung tepat.

\begin{proposisi}[Derau Gauss pada masukan sama dengan penalti L2]
\index{augmentasi data}
Misalkan $\hat y=\mathbf{w}^{\top}\mathbf{x}$ dan $\boldsymbol{\delta}\sim
N(\mathbf{0},\sigma^2I)$ ditambahkan pada masukan. Maka
\[
\mathbb{E}_{\boldsymbol{\delta}}\bigl[(\mathbf{w}^{\top}(\mathbf{x}+\boldsymbol{\delta})-y)^2\bigr]
=(\mathbf{w}^{\top}\mathbf{x}-y)^2+\sigma^2\|\mathbf{w}\|_2^2 .
\]
\end{proposisi}

\begin{proof}[Bukti]
Jabarkan: suku linearnya $2(\mathbf{w}^{\top}\mathbf{x}-y)\,
\mathbb{E}[\mathbf{w}^{\top}\boldsymbol{\delta}]=0$, dan suku kuadratnya
$\mathbb{E}[(\mathbf{w}^{\top}\boldsymbol{\delta})^2]
=\mathbf{w}^{\top}\mathbb{E}[\boldsymbol{\delta}\boldsymbol{\delta}^{\top}]\mathbf{w}
=\sigma^2\|\mathbf{w}\|_2^2$.
\end{proof}

Jadi menambahkan derau Gauss kecil pada masukan sama dengan menambahkan penalti L2
dengan $\lambda=\sigma^2$: keduanya satu gagasan yang ditulis dua kali. Pemeriksaan
angkanya, untuk $\mathbf{w}=(1,2)$, $\mathbf{x}=(1,0)$, $y=3$, dan $\sigma=0{,}1$:
sasaran asalnya $4$, ramalan identitasnya $4+0{,}01\cdot5=4{,}05$, dan simulasi satu juta
derau memberi $4{,}049107$. Selisih $0{,}0009$ itu sebanding galat acak simulasi, bukan
kesalahan rumus.

Peringatan untuk bidang keamanan: augmentasi hanya sah bila label benar-benar invarian
terhadap perubahan yang dilakukan. Menambahkan derau kecil pada ciri lalu lintas boleh,
tetapi membalik urutan paket atau menggeser waktu kedatangan dapat mengubah serangan
menjadi bukan serangan. Augmentasi yang salah label bukan regularisasi, melainkan
kebohongan yang akan tetap tampak baik selama data validasinya diubah dengan cara sama.

\section{SGD, momentum, dan Adam}
\tingkat{lanjut}

Sampai di sini kita memperlakukan "latihan" sebagai satu kotak. Isi kotak itu menentukan
seberapa cepat galat turun dan, pada anggaran langkah terbatas, seberapa rendah galat yang
tercapai. Kalimat sehari-harinya: arah yang benar belum cukup; yang menentukan adalah
bagaimana langkah diambil dan seberapa besar.

Penurun gradien stokastik mengganti gradien penuh oleh gradien atas satu minibatch
berukuran $b$. Taksiran itu tak bias, dan variansinya mengecil sebanding $1/b$:
$\operatorname{Var}(\hat{\mathbf{g}}_b)=\sigma_g^2/b$. Jadi minibatch besar mengurangi
goyangan, tetapi tidak mengubah arah rata-ratanya. Yang mengubah arah adalah momentum
dan Adam.

\begin{definisi}[Momentum dan Adam]
\index{momentum}
\index{adam}
Momentum menyimpan rata-rata gradien bergerak:
$\mathbf{v}_t=\beta\mathbf{v}_{t-1}+\mathbf{g}_t$ dan
$\mathbf{w}_t=\mathbf{w}_{t-1}-\eta\mathbf{v}_t$, dengan $0\le\beta<1$. Adam menyimpan
dua momen, $\mathbf{m}_t=\beta_1\mathbf{m}_{t-1}+(1-\beta_1)\mathbf{g}_t$ dan
$\mathbf{v}_t=\beta_2\mathbf{v}_{t-1}+(1-\beta_2)\mathbf{g}_t^2$ (kuadrat per
koordinat), lalu memakai koreksi bias dan pembaruan
$\mathbf{w}_t=\mathbf{w}_{t-1}-\eta\,\hat{\mathbf{m}}_t/(\sqrt{\hat{\mathbf{v}}_t}+\epsilon)$
dengan $\hat{\mathbf{m}}_t=\mathbf{m}_t/(1-\beta_1^t)$ dan
$\hat{\mathbf{v}}_t=\mathbf{v}_t/(1-\beta_2^t)$.
\end{definisi}

Koreksi bias itu bukan hiasan: tanpa koreksi $\mathbf{m}_1=(1-\beta_1)\mathbf{g}_1$
sangat kecil, sehingga langkah pertama hampir nol. Pembagian oleh $\sqrt{\hat{\mathbf{v}}}$
lebih penting lagi untuk masalah berskala buruk: ia membuat besar langkah tiap koordinat
kurang lebih $\eta$, tidak peduli gradiennya besar atau kecil. Momentum hanya meratakan
arah sepanjang waktu.

Angka pada kuadratik berikut memperlihatkan bedanya. Ambil
$f(\mathbf{w})=\frac12\mathbf{w}^{\top}H\mathbf{w}$ dengan $H=\operatorname{diag}(1,10^4)$
dan $\mathbf{w}_0=(1,0)$; bilangan kondisi $H$ adalah $10^4$, yang berarti kurvanya
sempit di satu arah dan sangat landai di arah lain. Kita hitung langkah sampai
$f\le10^{-8}f(\mathbf{w}_0)$.

\begin{center}
\small
\begin{tabular}{lrl}
\hline
metode & laju & langkah sampai $10^{-8}$ \\
\hline
penurun gradien & $10^{-4}$ & 92099 \\
momentum ($\beta=0{,}9$) & $10^{-4}$ & 9132 \\
Adam & $0{,}1$ & 82 \\
\hline
\end{tabular}
\end{center}

Penurun gradien memerlukan $92099$ langkah karena laju terbesarnya dibatasi kelengkungan
terbesar, sedangkan arah yang landai hanya bergerak dengan laju itu. Momentum memotong
satu digit, sesuai perkiraan bahwa percepatan mengubah ketergantungan pada bilangan
kondisi $\kappa$ menjadi sekitar $\sqrt\kappa$. Adam memotong tiga digit, sebab pembagian
per koordinat membuat arah yang landai tidak lagi dihukum oleh arah yang curam. Itu bukan
janji umum: pada masalah berskala baik keunggulan Adam menyusut.

Pada masalah logistik di bab ini, dengan minibatch delapan dan sasaran galat latih
$0{,}40$, median lima biji acak memberi SGD laju $0{,}05$ seratus langkah, sedangkan
Adam laju $0{,}1$ hanya sepuluh. Kesimpulannya harus hati-hati: keunggulannya sebanding
satu digit dan sebagian berasal dari laju yang lebih besar, bukan dari arah yang lebih
baik.

Laju pembelajaran biasanya tidak tetap. Dua jadwal yang lazim adalah peluruhan tangga,
yaitu laju dibagi sepuluh pada epoch tertentu, dan peluruhan kosinus
$\eta_t=\eta_{\min}+\frac12(\eta_0-\eta_{\min})(1+\cos(\pi t/T))$: langkah besar di awal
supaya cepat mendekat, langkah kecil di akhir supaya tidak melompat di lembah. Dengan
$3000$ langkah minibatch dan $\lambda=0{,}1$, laju tetap $0{,}2$ memberi galat latih
median $0{,}3866$, sedangkan kosinus dari $0{,}5$ ke $0{,}01$ memberi $0{,}3549$.

\section{Validasi silang, pemilihan hiperparameter, dan kebocoran}
\tingkat{lanjut}

Semua angka validasi di bab ini berasal dari himpunan yang tidak dipakai melatih. Itu
sudah benar, tetapi belum cukup: begitu kita memilih $\lambda$ dengan mata yang melihat
galat validasi, galat validasi itu berhenti menjadi taksiran yang tak bias bagi galat
populasi. Kalimat sehari-harinya: kalau kita memilih soal ujian berdasarkan nilai yang
kita dapat, nilai itu tidak lagi mengukur kemampuan kita.

\kamus{validasi silang}{cross validation}{memutar peran data bergiliran, sehingga setiap
contoh pernah menjadi ujian bagi model yang tidak melihatnya}

\begin{definisi}[Validasi silang $k$ lipatan]
\index{validasi silang}
Data latih dipecah menjadi $k$ lipatan yang saling lepas dan berukuran serupa. Untuk
setiap $j$, model dilatih pada $k-1$ lipatan dan diuji pada lipatan $j$; galat validasi
silang adalah rata-rata galat pada $k$ lipatan itu, dan hiperparameter yang dipilih
adalah yang meminimumkannya. Untuk $k=n$ prosedur ini disebut
\emph{leave-one-out}.
\end{definisi}

Validasi silang menjawab satu kelemahan pemisahan tunggal: pada $n=60$ himpunan
validasi berisi dua belas contoh, dan galat pada dua belas contoh bergoyang dengan
simpangan baku sekitar $0{,}12$. Dengan lima lipatan setiap contoh pernah diuji, tetapi
goyangannya tidak hilang.

\begin{center}
\small
\begin{tabular}{rrr}
\hline
$\lambda$ & galat CV (rata-rata 5 partisi) & simpangan antarpartisi \\
\hline
0     & 0,3233 & 0,1245 \\
0,003 & 0,2400 & 0,0253 \\
0,030 & 0,2467 & 0,0217 \\
0,100 & 0,2467 & 0,0274 \\
1     & 0,2967 & 0,0298 \\
\hline
\end{tabular}
\end{center}

Baris pertama menunjukkan dua hal sekaligus: tanpa regularisasi galat CV paling buruk
dan paling tidak stabil, dengan simpangan $0{,}1245$, empat kali lipat baris lainnya.
Baris kedua sampai keempat hampir tidak terbedakan: $0{,}2400$, $0{,}2467$, dan
$0{,}2467$, dengan simpangan sekitar $0{,}025$, dan pada lima partisi berbeda pemenangnya
berganti-ganti. Kurva CV ini datar di sekitar minimumnya, sehingga memilih nilai
persisnya adalah ilusi ketelitian.

Akibatnya nyata: CV memilih $\lambda=0{,}003$ dengan galat validasi $0{,}2688$,
sedangkan $\lambda=0{,}1$ yang tidak terpilih memberi $0{,}2406$; selisih $0{,}0282$ itu
tidak jauh dari satu simpangan baku kurva CV. Pelajarannya bukan bahwa validasi silang
rusak, melainkan bahwa dengan enam puluh contoh, pemilihan hiperparameter membawa
ketidakpastian sebesar efek yang diukur. Melaporkan "$\lambda$ terbaik $0{,}003$" tanpa
menyebut lebar kurvanya adalah kebohongan yang paling halus dalam bab
ini.

\begin{tangan}
Bahan: sepuluh kartu, pensil, kertas. Tulis satu bilangan dua angka pada tiap kartu,
kocok, lalu bentangkan tujuh kartu menghadap atas. Susun ``aturan'' dalam bentuk daftar:
posisi satu bernilai sekian, posisi dua bernilai sekian, sampai posisi tujuh.

Tutup satu kartu dengan telapak tangan dan tebak. Karena daftarmu memuat kartu itu,
tebakanmu tepat. Ulangi untuk ketujuh posisi: tujuh dari tujuh benar, dan itulah nilai
latihanmu.

Sekarang kumpulkan sepuluh kartu, kocok, bentangkan tujuh kartu baru, dan pakai daftar
yang sama. Tutup satu kartu, tebak, buka. Yang harus terlihat: daftarmu tidak memuat
kartu baru, jadi tebakanmu hanya untung-untungan. Ulangi sepuluh kali dengan susunan
baru dan catat berapa kali tepat; itulah nilai ujianmu. Terakhir tulis aturan kedua yang
memakai satu sifat saja, misalnya ``kartu di kiri tidak lebih besar dari kartu di
kanannya'', lalu bandingkan kedua nilai itu.
\end{tangan}

Ada sumber kesalahan yang lebih besar daripada goyangan itu dan tidak terlihat pada angka
mana pun: kebocoran, yaitu masuknya informasi data uji ke dalam proses latihan. Bentuk
paling umumnya adalah standardisasi ciri memakai rata-rata dan simpangan baku seluruh data
sebelum pemisahan dilakukan.

Angkanya, dengan pergeseran covariate buatan: bila validasi digeser $0{,}5$,
standardisasi dengan statistik latih saja memberi $0{,}3765$, sedangkan statistik
gabungan memberi $0{,}3361$; pada geseran $1{,}0$ angkanya $0{,}5387$ lawan $0{,}3819$.
Statistik gabungan selalu tampak lebih baik, dan justru itu masalahnya: ia memakai
informasi letak data uji yang tidak tersedia di lapangan. Tanpa pergeseran keduanya sama,
$0{,}2452$ lawan $0{,}2463$, tanda bahwa selisih tadi berasal dari kebocoran.

Protokol yang jujur berlapis: pisahkan latih, validasi, dan uji sejak awal; hitung
semua besaran, termasuk rata-rata, simpangan baku, ambang keputusan, dan pemilihan ciri,
di dalam lipatan latih saja; pilih hiperparameter pada validasi atau lewat validasi
silang di dalam data latih; sentuh himpunan uji satu kali, lalu laporkan apa adanya.

Hubungan kapasitas dengan jumlah data dapat dibuat kuantitatif. Untuk kelas hipotesis
berhingga, jurang generalisasi terkendali oleh akar logaritma cacah hipotesis dibagi
cacah contoh.

\begin{teorema}[Batas seragam untuk kelas berhingga]
\index{kapasitas}
Misalkan $\mathcal{H}$ berhingga dengan $|\mathcal{H}|=m$, galatnya 0-1, dan $n$ contoh
diambil bebas dari $\mathcal{D}$. Dengan peluang sekurang-kurangnya $1-\delta$ terhadap
pengambilan contoh, berlaku serentak untuk semua $h\in\mathcal{H}$
\[
\bigl|R(h)-\hat R_n(h)\bigr|\le\sqrt{\frac{\ln m+\ln(2/\delta)}{2n}} .
\]
\end{teorema}

\begin{proof}[Bukti]
Untuk satu $h$ tetap, $\hat R_n(h)$ adalah rata-rata $n$ peubah Bernoulli berskala
$[0,1]$, sehingga ketaksamaan Hoeffding memberi
$\mathbb{P}(|R(h)-\hat R_n(h)|>t)\le2e^{-2nt^2}$. Gabungkan atas semua $h$ dengan
ketaksamaan union: peluang ada satu $h$ yang melampaui $t$ paling besar
$2m\,e^{-2nt^2}$. Menetapkan ungkapan itu sama dengan $\delta$ memberi
$t=\sqrt{(\ln m+\ln(2/\delta))/(2n)}$, dan pelengkapnya adalah pernyataan teorema.
\end{proof}

Untuk kelas yang "menghafal", yaitu model yang memberi satu label pada setiap pola biner
atas $d$ ciri, $m=2^d$. Dengan $d=21$ dan $\delta=0{,}05$ batasnya
$\sqrt{(21\ln2+\ln40)/(2n)}$, yaitu $0{,}3899$ pada $n=60$ dan $0{,}0617$ pada
$n=2400$. Jurang terukurnya $0{,}1338$ dan $0{,}0097$ pada $n$ yang sama, jadi batasnya
berlaku tetapi longgar; kelonggaran itu bukan cacat, sebab batas itu harus berlaku untuk
setiap distribusi dan setiap algoritma.

Yang penting bukan angkanya, melainkan bentuknya. Jurang mengecil sebanding
$1/\sqrt n$, sedangkan cacah hipotesis masuk lewat logaritma. Melipatgandakan data empat
kali hanya memangkas jurang menjadi separuh, dan menambah kapasitas sepuluh kali lipat
hanya menaikkan batas sebesar akar penambahan logaritmanya. Inilah alasan jujur mengapa
model besar masih berguna: selama jurangnya dikendalikan regularisasi, kapasitas yang
naik tidak otomatis berarti jurang yang naik.

\section{Contoh hitung: ridge tertutup pada empat titik}
\tingkat{menengah}

Semua rumus di atas dapat diperiksa tangan pada contoh kecil: empat titik dengan satu
ciri, $\mathbf{x}=(1,2,3,4)$ dan $\mathbf{y}=(2,4,6,9)$, model tanpa intersep
$\hat y=wx$, dan sasaran $\frac12\sum_i(y_i-wx_i)^2+\frac{\lambda}{2}w^2$. Turunannya
nol memberi
\[
w(\lambda)=\frac{\sum_ix_iy_i}{\sum_ix_i^2+\lambda}=\frac{64}{30+\lambda},
\]
sebab $\sum x_i^2=30$ dan $\sum x_iy_i=2+8+18+36=64$.

\begin{center}
\small
\begin{tabular}{rrrr}
\hline
$\lambda$ & $w$ & SSE latih & $|w\cdot5-10|$ \\
\hline
0   & 2,133333 & 0,466667 & 0,666667 \\
1   & 2,064516 & 0,608741 & 0,322581 \\
2   & 2,000000 & 1,000000 & 0,000000 \\
5   & 1,828571 & 3,253061 & 0,857143 \\
10  & 1,600000 & 9,000000 & 2,000000 \\
30  & 1,066667 & 34,600000 & 4,666667 \\
100 & 0,492308 & 81,255621 & 7,538462 \\
\hline
\end{tabular}
\end{center}

Gambar~\ref{fig:lintasan-bobot} memplot kedua lintasan bobot terhadap $\lambda$ pada
data ini: L2 mengecil tanpa henti, L1 berhenti tepat di nol. Kolom terakhir memakai satu
titik uji di $x=5$ dengan nilai sebenarnya $10$, yang tidak ikut menentukan $w$. Di
situlah kurvanya berbentuk U dan minimumnya jatuh tepat di $\lambda=2$: di situ $w=2$,
sehingga ramalan di $x=5$ adalah $10$ dan galatnya nol. Tanpa regularisasi,
$w=64/30=2{,}1333$ lebih memikat pada data latih (SSE $0{,}466667$ lawan $1{,}000000$)
tetapi meramal $10{,}667$, meleset $0{,}667$. Inilah perdagangan bias-variansi dalam
bentuk paling telanjang: SSE latih naik dua kali, galat uji menjadi nol.

\begin{figure}[htbp]
\centering
\begin{tikzpicture}
\begin{axis}[
  width=11.8cm, height=5.4cm,
  xmode=log,
  xmin=0.5, xmax=300, ymin=-0.12, ymax=2.45,
  axis lines=left,
  xtick={1,2,10,64,100},
  xticklabels={$1$,$2$,$10$,$64$,$100$},
  ytick={0,0.5,1.0,1.5,2.0},
  yticklabels={$0$,$0{,}5$,$1{,}0$,$1{,}5$,$2{,}0$},
  tick label style={font=\scriptsize},
  xlabel={kekuatan regularisasi $\lambda$ (skala log)},
  ylabel={bobot $w$},
  xlabel style={font=\small},
  ylabel style={font=\small},
  legend style={font=\scriptsize,draw=none,at={(0.97,0.97)},anchor=north east},
]
\addplot[domain=0.5:300,samples=200,very thick,layerindigo] {64/(30+x)};
\addlegendentry{L2: $64/(30+\lambda)$}
\addplot[domain=0.5:63.9,samples=200,very thick,layerorange] {(64-x)/30};
\addlegendentry{L1: $\max\{(64-\lambda)/30,\,0\}$}
\addplot[domain=64:300,samples=2,very thick,layerorange] {0};
\draw[densely dashed,layergrey] (axis cs:2,-0.12)--(axis cs:2,2.20);
\draw[densely dashed,layergrey] (axis cs:64,-0.12)--(axis cs:64,2.20);
\addplot[only marks,mark=*,mark size=1.8pt,layerindigo] coordinates {(2,2.0000) (64,0.6809)};
\addplot[only marks,mark=square*,mark size=1.8pt,layerorange] coordinates {(2,2.0667)};
\addplot[only marks,mark=o,mark size=3.4pt,layerorange] coordinates {(64,0)};
\node[font=\scriptsize,anchor=south] at (axis cs:2,2.24) {$\lambda=2$};
\node[font=\scriptsize,anchor=south] at (axis cs:64,2.24) {$\lambda=64$};
\node[font=\scriptsize,anchor=south west] at (axis cs:70,0.06) {L1 tepat nol};
\end{axis}
\end{tikzpicture}
\caption{Lintasan bobot terhadap kekuatan regularisasi pada contoh hitung empat titik,
dengan $\sum x_i^2=30$ dan $\sum x_iy_i=64$. Penalti L2 memberi
$w=64/(30+\lambda)$, yang mengecil tetapi tidak pernah menyentuh nol. Penalti L1 memberi
$w=\max\{(64-\lambda)/30,\,0\}$, yang mencapai nol tepat pada $\lambda=64$. Pada
$\lambda=2$ keduanya hampir sama, $2{,}0000$ lawan $2{,}0667$; pada $\lambda=64$ L2 masih
$0{,}6809$ sedangkan L1 sudah $0{,}0000$. Garis putus-putus menandai kedua nilai
$\lambda$ yang dibandingkan.}
\label{fig:lintasan-bobot}
\end{figure}

Perhatikan dinding kanannya. Pada $\lambda=100$, $w$ tinggal $0{,}4923$, SSE latih
meledak menjadi $81{,}255621$, dan galat uji $7{,}538462$, lebih buruk daripada model
yang tidak belajar apa pun. Contoh ini juga menunjukkan mengapa mencari $\lambda$ pada
data latih menyesatkan: kurva SSE latih monoton naik, jadi ia selalu memilih
$\lambda=0$.

\begin{tantangan}
Sebuah model melaporkan galat latih $0{,}0000$ pada $60$ contoh dan galat uji
$0{,}4100$. Rekannya mengusulkan tiga perubahan: (i) melatih lebih lama, (ii) menambah
ciri sampai derajat tujuh, (iii) menaikkan $\lambda$ dari $0$ ke $0{,}1$. Untuk setiap
usulan, ramalkan tanda perubahan galat latih dan galat uji, lalu sebutkan mana yang
tidak mungkin menolong beserta alasannya.

Petunjuk 1: tanyakan lebih dulu ke arah mana kapasitas bergerak pada tiap usulan.

Petunjuk 2: dua dari tiga usulan memperbesar kapasitas, dan model ini sudah berada di
galat latih nol, sehingga memperbesar kapasitas hanya memperlebar jurang. Yang menolong
adalah yang mengecilkannya, dan itu harus menaikkan galat latih.

Kalau mentok, kembali ke Definisi underfitting dan overfitting pada bagian pertama.
\end{tantangan}

\kertas{Periksa kolom kedua tanpa kalkulator: $64/32=2$ tepat, $64/31=2{,}064516$, dan
$64/40=1{,}6$. Lalu hitung sendiri SSE pada $\lambda=2$: sisaannya $(0,0,0,1)$, jadi
SSE $=1$. Terakhir jawab: berapa $\lambda$ yang membuat $w=1{,}28$, dan berapa ramalan
di $x=5$ ketika itu? Petunjuk: selesaikan $64/(30+\lambda)=1{,}28$.}

\section{Lab komputasi}
\tingkat{lanjut}

Seluruh angka pada bab ini berasal dari satu skrip numpy tanpa pustaka pembelajaran
mesin. Datanya sengaja kecil supaya dapat diperiksa ulang dalam hitungan menit: dua ciri
mentah, dua puluh satu ciri polinomial sampai derajat lima, enam puluh contoh latih, dan
dua puluh ribu contoh validasi dari distribusi yang sama. Labelnya tidak deterministik,
sehingga ada galat tak tereduksikan sekitar $0{,}21$ yang tidak dapat dilampaui model apa
pun.

\begin{lab}
\begin{verbatim}
import numpy as np

def data(n, rng):
    x = rng.uniform(-1, 1, size=(n, 2))
    p = 1/(1 + np.exp(-4.0*(0.9 - x[:,0]**2 - x[:,1]**2)))
    return x, (rng.random(n) < p).astype(float)

def poly(x, deg=5):
    return np.column_stack([x[:,0]**i * x[:,1]**j
                            for i in range(deg+1) for j in range(deg+1-i)])

def sig(z):
    return 1/(1 + np.exp(-np.clip(z, -400, 400)))

def soft(t, l):
    return np.sign(t)*np.maximum(np.abs(t) - l, 0.0)

def er(Z, y, w, b):
    return float(np.mean((sig(Z@w + b) > 0.5) != y))

def fit(Z, y, lam, iters=30000, lr=0.05):
    n, d = Z.shape
    w = np.zeros(d); b = 0.0
    mw = np.zeros(d); vw = np.zeros(d); mb = 0.0; vb = 0.0
    for t in range(1, iters+1):
        p = sig(Z@w + b); g = (p - y)/n
        gw = Z.T@g + lam*w; gb = g.sum()
        mw = 0.9*mw + 0.1*gw; vw = 0.999*vw + 0.001*gw*gw
        mb = 0.9*mb + 0.1*gb; vb = 0.999*vb + 0.001*gb*gb
        w -= lr*(mw/(1-0.9**t))/(np.sqrt(vw/(1-0.999**t)) + 1e-8)
        b -= lr*(mb/(1-0.9**t))/(np.sqrt(vb/(1-0.999**t)) + 1e-8)
    return w, b

def fit_l1(Z, y, lam, iters=150000):
    n, d = Z.shape
    w = np.zeros(d); b = 0.0
    st = 0.5/(0.25*np.linalg.eigvalsh(Z.T@Z/n).max())
    for _ in range(iters):
        p = sig(Z@w + b); g = (p - y)/n
        w = soft(w - st*(Z.T@g + lam*w), st*lam); b -= st*g.sum()
    return w, b
\end{verbatim}
\end{lab}

Bagian pertama menghitung galat untuk sebelas nilai $\lambda$; bagian kedua mengulanginya
dengan penalti L1 lewat penurunan gradien proksimal, yaitu langkah gradien yang diikuti
ambang lunak dari lemma tadi; bagian ketiga berhenti lebih awal; bagian keempat dan
kelima membandingkan pengoptimum; bagian keenam sampai kedelapan memeriksa validasi
silang, kebocoran, dan kapasitas. Keluaran sebenarnya, dipotong hanya pada bagian yang
berulang:

\begin{verbatim}
d=21  n_latih=60  n_validasi=20000
==1 L2: lam  |w|max  latih  validasi
  0.0000 406.204  0.0000  0.3262
  0.0010   3.056  0.1000  0.2927
  0.0030   1.814  0.1000  0.2688
  0.0100   1.018  0.1167  0.2475
  0.0300   0.631  0.1167  0.2433
  0.1000   0.399  0.1500  0.2406
  0.3000   0.268  0.2333  0.2452
  1.0000   0.142  0.2500  0.2695
  3.0000   0.062  0.3333  0.3205
 10.0000   0.022  0.3333  0.3395
100.0000   0.003  0.3333  0.3395
==2 L1: lam  nol/total  latih  validasi
  0.0020 16/21  0.1167  0.2630
  0.0050 11/21  0.1167  0.2455
  0.0100 14/21  0.1167  0.2412
  0.0200 14/21  0.1500  0.2369
  0.0500 11/21  0.1833  0.2334
  0.1000  7/21  0.2667  0.2492
  0.2000  1/21  0.3333  0.3395
==3 early stopping: epoch  latih  validasi  |w|max
ep=   1  0.1333  0.2481  0.15
ep=   5  0.1333  0.2429  0.38
ep=  10  0.1333  0.2432  0.48
ep=  50  0.1167  0.2454  0.77
ep= 200  0.1167  0.2541  1.38
ep= 800  0.1000  0.2879  2.85
ep=2400  0.0833  0.3109  4.81
terbaik 0.2427 pada epoch 8
==4 kuadratik diag(1,1e4): langkah ke galat 1e-8
GD=1e-4  92099 | momentum  9132 | Adam=0.1  82
==5 logistik minibatch: langkah ke galat latih 0.40 (median 5 seed)
sgd  lr=0.05 [100, 100, 90, 90, 100] median=100
sgd  lr=0.2  [30, 50, 40, 40, 20] median=40
adam lr=0.03 [30, 20, 20, 20, 30] median=20
adam lr=0.1  [20, 10, 20, 10, 10] median=10
==6 CV 5 partisi x 5 lipatan: lam  galat CV (rata2 +- sd antarpartisi)
 0.000  0.3233 +- 0.1245
 0.003  0.2400 +- 0.0253
 0.030  0.2467 +- 0.0217
 0.100  0.2467 +- 0.0274
 1.000  0.2967 +- 0.0298
==7 standardisasi: geser  statistik latih  statistik gabungan
0.0  0.2452  0.2463
0.5  0.3765  0.3361
1.0  0.5387  0.3819
==8 kapasitas vs data (d=21, lam=0, himpunan bersarang)
n=   60 latih=0.1667 validasi=0.3005 selisih=+0.1338 batas=0.3899
n=  150 latih=0.2133 validasi=0.2407 selisih=+0.0274 batas=0.2466
n=  600 latih=0.2067 validasi=0.2260 selisih=+0.0193 batas=0.1233
n= 2400 latih=0.2058 validasi=0.2155 selisih=+0.0097 batas=0.0617
\end{verbatim}

Tiga hal perlu dibaca dari keluaran itu, dan ketiganya mengecewakan harapan biasa.

Pertama, galat latih hampir tidak berguna sebagai petunjuk: ia bergerak berlawanan arah
dengan galat validasi. Kurva validasinya juga tidak simetris. Dari $\lambda=0$ ke
$0{,}1$ galatnya turun $0{,}086$; dari $0{,}1$ ke $10$ ia naik $0{,}099$ sampai melewati
titik awalnya. Bila ragu antara dua nilai, yang lebih kecil lebih aman.

Kedua, pemilihan lewat validasi silang tidak otomatis benar. CV memilih
$\lambda=0{,}003$, sedangkan validasi memilih $0{,}1$; selisih galat CV keduanya
$0{,}0006$, jauh di bawah simpangan antarpartisi $0{,}025$, padahal di data validasi
keduanya berbeda $0{,}0282$.

Ketiga, batas teoretis longgar tetapi arahnya benar. Eksperimen kapasitas di sini
melatih model dengan anggaran enam puluh ribu langkah pada himpunan bersarang, jadi galat
latihnya tidak nol. Pada $n=60$ batasnya $0{,}3899$ sedangkan jurangnya $0{,}1338$; pada
$n=2400$ batasnya $0{,}0617$ sedangkan jurangnya $0{,}0097$. Rasionya sekitar tiga,
sehingga batas itu berguna untuk membandingkan setelan, bukan meramal angka.

\begin{untukdosen}
Urutan mengajar yang saya sarankan: mulai dari angka pembuka bab, enam puluh contoh
dengan galat latih nol dan galat validasi $0{,}3262$, lalu tanyakan apakah model itu
"belajar". Jawaban kelas biasanya membelah dua, dan perpecahan itu bahan kuliah terbaik.
Sesudah itu definisi formal, dan biarkan dekomposisi bias-variansi diturunkan di papan
dalam lima baris. Regularisasi L1 sebaiknya masuk lewat lemma ambang lunak, bukan lewat
gambar bola L1, sebab lemma itu dapat diperiksa mahasiswa sendiri.

Dua miskonsepsi yang paling sering muncul. Mahasiswa menganggap galat latih rendah sebagai
bukti model bagus, lalu terkejut ketika galat validasinya tinggi; gejala lain, mereka
menyarankan "melatih lebih lama" sebagai obat, padahal pada tabel berhenti lebih awal
justru itulah yang merusak. Dan mereka menganggap regularisasi selalu menolong, lalu
memasang $\lambda$ besar dengan keyakinan "lebih aman". Tabel $\lambda=10$ dengan galat
$0{,}3395$ adalah obat untuk keduanya.

Aktivitas papan tulis yang berhasil: minta kelas menebak, sebelum dihitung, di mana
minimum kurva validasi berada, lalu hitung bersama tiga baris saja; tebakan biasanya
meleset ke $\lambda$ yang terlalu kecil. Tanda kelas siap lanjut: mereka dapat menjelaskan
mengapa memilih $\lambda$ memakai data validasi membuat angka itu optimistis, dan dapat
menyebutkan satu kebocoran yang tidak terlihat pada angka.
\end{untukdosen}

\section{Ringkasan}
\tingkat{dasar}

\begin{ringkasan}
\begin{itemize}
  \item Galat latih dan galat populasi dihitung dengan cara sama pada himpunan berbeda;
        galat latih $0{,}0000$ dengan galat validasi $0{,}3262$ menandakan hafalan.
  \item Galat populasi terdekomposisi menjadi derau tak tereduksi, bias kuadrat, dan
        variansi; regularisasi menukar bias dengan variansi, dan perdagangan itu
        menguntungkan hanya selama penurunan variansi lebih besar daripada kenaikan bias.
  \item Penalti L2 menambah $\lambda\mathbf{w}$ pada gradien, setara mengalikan bobot
        dengan $1-\eta\lambda$ tiap langkah; L1 memakai ambang lunak sehingga bobot di
        bawah $\lambda/H$ menjadi nol tepat. Keduanya punya dosis: $\lambda=0{,}1$
        memberi $0{,}2406$, sedangkan $\lambda=10$ memberi $0{,}3395$, lebih buruk
        daripada tanpa regularisasi.
  \item Berhenti lebih awal adalah regularisasi lewat jumlah langkah: $0{,}2427$ pada
        epoch 8 lawan $0{,}3109$ pada epoch 2400. Dropout menambahkan derau bervariansi
        $(1-p)/p$ per unit, dan augmentasi derau Gauss $\sigma^2$ setara penalti L2
        berskala $\sigma^2$.
  \item Pengoptimum berbeda pada masalah berskala buruk: pada kuadratik berbilangan
        kondisi $10^4$, penurun gradien perlu $92099$ langkah, momentum $9132$, Adam
        $82$; jadwal kosinus menurunkan galat latih dari $0{,}3866$ ke $0{,}3549$.
  \item Validasi silang mengurangi goyangan tanpa menghapusnya: kurva CV datar dengan
        simpangan $0{,}025$, dan pilihannya meleset $0{,}0282$. Kebocoran standardisasi
        membuat angka tampak lebih baik ($0{,}3361$ lawan $0{,}3765$) tanpa benar-benar
        lebih baik, dan jurang generalisasi menyusut sebanding $1/\sqrt n$: dari
        $0{,}1338$ pada $n=60$ ke $0{,}0097$ pada $n=2400$.
\end{itemize}
\end{ringkasan}

\section{Latihan}
\tingkat{menengah}

\begin{latihan}
\begin{enumerate}
  \item Nilai klaim ini: "Karena galat latih model kami nol, model itu sudah cukup baik."
        Benar, salah, atau belum cukup bukti? Sebutkan bagian yang keliru dan satu angka
        tambahan yang Anda perlukan.
  \item Buktikan dekomposisi bias-variansi untuk $y=f(x)+\varepsilon$ dengan
        $\operatorname{Var}(\varepsilon)=\sigma^2$: tuliskan ketiga suku, tunjukkan kedua
        suku silang lenyap, dan sebutkan di mana kebebasan $D$ dari $\varepsilon$
        dipakai.
  \item Turunkan $w(\lambda)=64/(30+\lambda)$ untuk data pada contoh hitung, lalu
        buktikan $w$ turun monoton terhadap $\lambda$ tetapi tidak pernah nol. Bandingkan
        dengan penalti L1 pada masalah satu dimensi yang sama, dan tentukan $\lambda$
        yang membuat $w=0$ tepat.
  \item Hitung galat validasi untuk $c=0{,}95$ dan $c=1{,}05$ pada contoh penaksir
        menyusut dengan $\mu=2$, $\sigma=1$, $n=25$, lalu bandingkan dengan $c=1$.
        Jelaskan mengapa dua simpangan sama besar tidak memberi kerugian sama besar.
  \item Sebuah laporan menulis: "Kami memilih $\lambda$ terbaik lewat validasi silang,
        jadi angka validasi silang itu adalah perkiraan galat di lapangan." Perbaiki
        kalimat itu, dan sebutkan dua prosedur tambahan yang diperlukan.
  \item Buktikan lemma ambang lunak untuk $a<0$ secara lengkap, lalu jelaskan mengapa
        hasilnya simetris terhadap tanda $a$. Petunjuk: kerjakan kasus $u<0$ dan $u=0$
        terpisah.
  \item Hitung langkah peluruhan bobot yang diperlukan agar bobot menjadi seperempat
        nilai awalnya bila $\eta=0{,}1$ dan $\lambda=0{,}02$, memakai waktu paruh
        $\ln2/(\eta\lambda)$. Periksa dengan simulasi numpy untuk gradien nol.
  \item Pada data bab ini, $\lambda=0{,}1$ memberi galat latih $0{,}1500$ dan galat
        validasi $0{,}2406$. Hitung jurangnya, bandingkan dengan jurang pada $\lambda=0$
        dan $\lambda=10$, lalu jelaskan mengapa jurang kecil tidak cukup untuk
        menyimpulkan sebuah model lebih baik.
  \item Turunkan batas seragam $\sqrt{(\ln m+\ln(2/\delta))/(2n)}$ dari ketaksamaan
        Hoeffding dan union bound, lalu hitung untuk $d=10$, $n=200$, dan
        $\delta=0{,}01$. Berapa kali data harus dilipatgandakan agar batasnya separuh?
\end{enumerate}

Petunjuk untuk tiga butir tersulit. Butir 3: sasaran satu dimensinya
$\frac12\sum_i(y_i-wx_i)^2+\lambda|w|$; pada $w=0$ subgradiennya memuat nol bila
$\lambda\ge|\sum_ix_iy_i|=64$, jadi $w=0$ tepat pada $\lambda=64$. Butir 6: pada $u=0$
subgradiennya selang $[-a-t,\,-a+t]$, yang memuat nol bila $|a|\le t$. Butir 9: dengan
$m=2^{10}$ dan $\delta=0{,}01$ batasnya $\sqrt{(10\ln2+\ln200)/400}=0{,}1749$;
separuhnya menuntut empat kali data.
\end{latihan}

\begin{untukpeneliti}
Yang terbukti di bab ini: dekomposisi bias-variansi beserta syaratnya; kesetaraan
langkah peluruhan bobot dengan penalti L2; bentuk tertutup peminimum
$\frac H2(w-w^\ast)^2+\lambda|w|$ dan syarat subgradien yang membuat L1 menghasilkan nol
sejati; kesetaraan augmentasi derau Gauss dengan penalti L2 untuk model linear; rumus
harapan dan variansi topeng dropout; serta batas seragam untuk kelas hipotesis
berhingga. Yang hanya populer: bahwa regularisasi selalu menolong; bahwa Adam selalu
lebih cepat; bahwa berhenti lebih awal "sama dengan" L2 (itu aproksimasi untuk sasaran
kuadratik dengan laju kecil); bahwa cacah nol L1 mengukur kepentingan ciri.

Sumber primer: Goodfellow, Bengio, dan Courville, \emph{Deep Learning}, Bab 5
(kapasitas, bias-variansi), Bab 7 (penalti L1 dan L2, augmentasi, derau, berhenti lebih
awal, dropout), Bab 8 (SGD, momentum, Adam, jadwal laju, pemilihan hiperparameter).
Wasserman, \emph{All of Statistics}, Bab 13 dan Bab 20. Batas seragam adalah akibat
langsung ketaksamaan Hoeffding dan union bound.

Batas keberlakuan: seluruh angka dari satu himpunan sintetis dengan pemisahan tunggal
$60/20000$, sehingga tiap galat validasi membawa ketidakpastian orde $0{,}01$ sampai
$0{,}02$. Karena itu perbandingan L1 dan L2 di sini bukan rekomendasi. Batas seragam hanya
untuk kelas berhingga dan galat 0-1.

Pertanyaan terbuka: seberapa besar kebocoran normalisasi batch, yang menghitung
statistik dari minibatch sehingga mencampur informasi antarcontoh; bagaimana memilih
$\lambda$ ketika biaya kesalahan tidak simetris; dan apakah kurva CV yang datar dapat
diubah menjadi interval yang sah, bukan sekadar simpangan antarpartisi.
\end{untukpeneliti}

\peta{Bab 26 pada Goodfellow, Bengio, dan Courville, \emph{Deep Learning}: Bab 5
(kapasitas, overfitting dan underfitting, dekomposisi bias-variansi), Bab 7 (penalti L1
dan L2, weight decay, augmentasi data, derau sebagai regularisasi, berhenti lebih awal,
dropout), Bab 8 (SGD, momentum, Adam, jadwal laju, pemilihan hiperparameter). Pada
Wasserman, \emph{All of Statistics}: Bab 13 (pemilihan model dan validasi silang), Bab 20
(bias-variansi pada penghalusan). Turunan balik dan pelatihan jaringan dibahas di Bab 23;
pembacaan metrik evaluasi yang jujur, termasuk kebocoran dan pemisahan data, dilanjutkan
di Bab 25; optimasi konveks yang mendasari penalti dan subgradien ada di Bab 21.}
